\documentclass[11pt,letterpaper]{article}
\usepackage{fontspec}
\usepackage{tocloft}
\usepackage{multicol}
\usepackage{fancyhdr}
\usepackage{nicefrac}
\usepackage[left=1.3in,right=1.3in,top=1in,bottom=1in]{geometry}
\setmainfont[Scale=1.3,AutoFakeBold=1.5,AutoFakeSlant=0.3]{Doves Type}

\widowpenalty=10000
\clubpenalty=10000
\interlinepenalty=500

\title{Lamb Birria Tacos}
\author{}
\date{}

% Configure fancy header
\pagestyle{fancy}
\fancyhf{} % Clear all header and footer fields
\fancyhead[C]{\textit{Lamb Birria Tacos}} % Center header with recipe title
\fancyfoot[C]{\thepage} % Center page number at bottom
\renewcommand{\headrulewidth}{0pt} % Remove header line
\renewcommand{\footrulewidth}{0pt} % Remove footer line

\begin{document}

\maketitle
\thispagestyle{empty}

\begin{quote}\textit{A bone-in lamb leg is roasted, then braised in a brick-red adobo of guajillo, ancho, and árbol chiles warmed with cumin, clove, and cinnamon until the meat falls from the bone. The shredded lamb is moistened with its own consommé, then griddled into corn tortillas with melted queso Oaxaca to make crisp quesabirria, served with the reserved consommé for dipping. Serves 6--8 (about 16--20 tacos).}\end{quote}

\section*{Ingredients}
\setlength{\columnsep}{20pt}
\begin{multicols}{2}
\noindent
    Bone-in lamb leg \dotfill 4 lbs. \\
    Kosher salt and black pepper \dotfill as needed \\
    Guajillo chiles \dotfill 12 \\
    Ancho chiles \dotfill 5 \\
    Árbol chiles \dotfill 3 \\
    Water, for soaking chiles \dotfill 2 qts. \\
    Tomatoes \dotfill 3 \\
    Garlic cloves, crushed \dotfill 6 \\
    White onion, chopped \dotfill 1 \\
    Ground cumin \dotfill 2 tsp. \\
    Mexican oregano \dotfill 2 tsp. \\
    \columnbreak
    Black peppercorns \dotfill 1 tsp. \\
    Whole cloves \dotfill \nicefrac{3}{4} tsp. \\
    Ground cinnamon \dotfill \nicefrac{1}{4} tsp. \\
    Kosher salt \dotfill 4 Tbsp. \\
    White vinegar \dotfill \nicefrac{1}{3} cup \\
    Vegetable oil \dotfill \nicefrac{1}{3} cup \\
    Bay leaf \dotfill 1 \\
    Corn tortillas \dotfill 16--20 \\
    Queso Oaxaca, shredded \dotfill 2--3 cups \\
    White onion, finely chopped \dotfill 1 cup \\
    Fresh cilantro, chopped \dotfill 1 cup \\
    Lime wedges \dotfill for serving \\
\end{multicols}

\filbreak
\section*{The Lamb}

\noindent
Preheat oven to \textit{325°F} ---
Trim silver skin and excess fat from \textbf{lamb~leg}; pat dry

\begin{enumerate}
    \item Cut 3--4 deep scores into the \textbf{lamb~leg} and season it all over, generously, with \textbf{salt} and \textbf{black~pepper}.
    \item Wrap the \textbf{lamb} loosely in foil, set it on a roasting tray, and roast at \textit{325°F} for \textit{2\nicefrac{1}{2}~hours}. Open the foil for the final \textit{30~minutes} so the top browns. Make the adobo while the \textbf{lamb} roasts.
\end{enumerate}

\filbreak
\section*{Birria Adobo}

\noindent
Stem, seed, and rinse \textbf{guajillo}, \textbf{ancho}, and \textbf{árbol~chiles} ---
Crush \textbf{garlic} and roughly chop \textbf{onion}; combine in \textit{Medium Bowl~\#1} (aromatics) ---
Combine 2~tsp. \textbf{cumin}, 2~tsp. \textbf{oregano}, 1~tsp. \textbf{peppercorns}, \nicefrac{3}{4}~tsp. \textbf{cloves}, and \nicefrac{1}{4}~tsp. \textbf{cinnamon} in \textit{Small Bowl~\#1} (spices) ---
Measure 4~Tbsp. \textbf{salt} and \nicefrac{1}{3}~cup \textbf{vinegar}

\begin{enumerate}
    \item Put the \textbf{chiles} in a large pot with 2~qts. cold \textbf{water}. Bring to a boil, covered, then turn off the heat and leave covered for \textit{10~minutes}, until soft and pliable.
    \item Meanwhile, boil the \textbf{tomatoes} in a separate saucepan of water for \textit{3~minutes}, until the skins loosen and split. Drain and set aside in \textit{Medium Bowl~\#2}.
    \item Drain the \textbf{chiles}, saving the soaking liquid in \textit{Large Bowl~\#1}---it carries flavor and adjusts the consommé later. Move the \textbf{chiles} to \textit{Large Bowl~\#2}.
    \item In a blender, combine the \textbf{chiles} (\textit{Large Bowl~\#2}), \textbf{tomatoes} (\textit{Medium Bowl~\#2}), \textbf{aromatics} (\textit{Medium Bowl~\#1}), \textbf{spices} (\textit{Small Bowl~\#1}), 4~Tbsp. \textbf{salt}, \nicefrac{1}{3}~cup \textbf{vinegar}, and 4~cups reserved \textbf{chile~water} (\textit{Large Bowl~\#1}). Blend smooth, about \textit{1--2~minutes}, in batches if the jar is full.
    \item Heat \nicefrac{1}{3}~cup \textbf{oil} in a large pot over medium heat until shimmering. Add the blended adobo---stand back, it splatters---and the \textbf{bay~leaf}. Cook, stirring often, for \textit{20~minutes}, until it deepens in color and thickens slightly and a layer of red \textbf{chile~fat} rises to the top. Taste, adjust \textbf{salt}, and turn off the heat.
\end{enumerate}

\filbreak
\section*{Braise and Consommé}

\begin{enumerate}
    \item Once the \textbf{lamb} has roasted \textit{2\nicefrac{1}{2}~hours}, remove the foil and pour in as much adobo as the tray holds, letting it mingle with the \textbf{lamb} juices; reserve the rest in the pot. Cover tightly with foil and braise at \textit{325°F} for \textit{2 to 2\nicefrac{1}{2}~hours}, flipping the \textbf{lamb} halfway, until it is fork-tender and the bone slips out cleanly (\textit{195--205°F}).
    \item Just before the \textbf{lamb} is done, skim the red \textbf{chile~fat} off the reserved adobo into \textit{Small Bowl~\#2}; this is the frying fat for the tortillas.
    \item Drain the tray juices back into the reserved adobo and stir. Loosen to a sippable consommé with reserved \textbf{chile~water}, starting with 2~cups and adding more to taste, then strain through a fine-mesh sieve and keep warm.
    \item Rest the \textbf{lamb} \textit{30~minutes}, then remove the bone and shred the meat into \textit{Large Bowl~\#3}, discarding fat and gristle. Toss with about 1~cup of the warm consommé to moisten.
\end{enumerate}

\filbreak
\section*{Quesabirria Tacos}

\noindent
Finely chop garnish \textbf{onion} and \textbf{cilantro}; combine in \textit{Small Bowl~\#3} (garnish) ---
Cut \textbf{lime} into wedges

\begin{enumerate}
    \item Heat a large skillet or comal over medium-high heat until hot but not smoking. Dip a \textbf{tortilla} in the \textbf{chile~fat} (\textit{Small Bowl~\#3}), coating both sides, and lay it in the pan.
    \item Scatter on \textbf{queso~Oaxaca}, then a portion of \textbf{lamb}, then a spoonful of consommé. Fold the \textbf{tortilla} over with a spatula and press gently.
    \item Cook for \textit{2--3~minutes}, until the underside is crisp and deep red, then flip and cook \textit{2--3~minutes} more, until both sides are crisp and the cheese melts. Repeat with the remaining \textbf{tortillas}.
    \item Garnish with the \textbf{onion-cilantro} mixture (\textit{Small Bowl~\#3}) and a squeeze of \textbf{lime}. Serve right away with bowls of warm consommé for dipping.
\end{enumerate}

\filbreak

\newpage

% Begin compact two-column layout
{\small
\setlength{\columnsep}{20pt}
\setlength{\multicolsep}{6pt}
\begin{multicols}{2}
\setlength{\parindent}{0pt}
\setlength{\parskip}{4pt}

\section*{Equipment Required}
\begin{itemize}
    \item Roasting tray and heavy-duty aluminum foil
    \item Large pot (at least 4-quart) for soaking chiles
    \item Small saucepan for the tomatoes
    \item High-powered blender (at least 500 watts)
    \item Large pot or Dutch oven (at least 6-quart) for frying the adobo
    \item Fine-mesh sieve or strainer
    \item Large (12-inch) skillet, comal, or griddle for the tacos
    \item Metal spatula and kitchen tongs
    \item Two forks for shredding
    \item Mixing bowls (small, medium, large)
    \item Sharp chef's knife and kitchen shears
    \item Cutting board
    \item Instant-read thermometer
    \item Ladle and small bowls for serving consommé
\end{itemize}

\section*{Hints and Notes}

\subsection*{Yield}
\begin{itemize}
    \item Serves 6--8, making about 16--20 tacos with consommé for dipping.
\end{itemize}

\subsection*{Mise en Place}
\begin{itemize}
    \item \textbf{Small Bowl~\#1} --- spices: 2~tsp. \textbf{cumin}, 2~tsp. \textbf{oregano}, 1~tsp. \textbf{peppercorns}, \nicefrac{3}{4}~tsp. \textbf{cloves}, \nicefrac{1}{4}~tsp. \textbf{cinnamon}
    \item \textbf{Small Bowl~\#2} --- reserved \textbf{chile~fat} skimmed from the adobo (for frying tortillas)
    \item \textbf{Small Bowl~\#3} --- garnish: 1~cup finely chopped \textbf{onion} and 1~cup chopped \textbf{cilantro}
    \item \textbf{Medium Bowl~\#1} --- aromatics: crushed \textbf{garlic} and chopped \textbf{onion}
    \item \textbf{Medium Bowl~\#2} --- boiled \textbf{tomatoes}, set aside
    \item \textbf{Large Bowl~\#1} --- reserved \textbf{chile~water} (about 6~cups; 4~cups for the blender, the rest to adjust the consommé)
    \item \textbf{Large Bowl~\#2} --- drained, softened \textbf{chiles}
    \item \textbf{Large Bowl~\#3} --- shredded \textbf{lamb} tossed with consommé
    \item Measure 4~Tbsp. \textbf{salt} and \nicefrac{1}{3}~cup \textbf{vinegar} for the blender; keep the \textbf{bay~leaf} and \nicefrac{1}{3}~cup \textbf{oil} by the stove
    \item Take the \textbf{lamb} out \textit{30--45~minutes} ahead so it roasts evenly
\end{itemize}

\subsection*{Ingredient Tips}
\begin{itemize}
    \item \textbf{Lamb~leg}: bone-in delivers the most flavor and gelatin, and the bone slides free when the meat is done. A bone-in shoulder works as a richer, fattier substitute. Have the butcher leave the shank on if it fits your tray.
    \item \textbf{Chiles}: choose pliable, leathery pods that bend without cracking---brittle, dusty \textbf{chiles} are stale and turn bitter. \textbf{Guajillo} brings bright color, \textbf{ancho} sweet depth, and \textbf{árbol} the heat; pull one or two \textbf{árbol} for a milder batch.
    \item \textbf{Queso~Oaxaca}: the traditional melting cheese for quesabirria---mild and stringy, it pulls into threads on the griddle. Asadero or Chihuahua melt similarly; mozzarella is an acceptable last resort.
    \item \textbf{Mexican~oregano}: more citrusy and less minty than Mediterranean oregano, and the authentic choice here. If unavailable, use the same amount of Mediterranean oregano and accept a slightly sweeter herb note.
    \item \textbf{Corn~tortillas}: thicker, fresh tortillas hold up best to dipping and griddling. Warm them briefly if they crack when folded.
\end{itemize}

\subsection*{Preparation Tips}
\begin{itemize}
    \item \textbf{Frying~the~adobo}: this step is not optional---raw blended chile tastes flat, and the tacos rely on the rendered fat. Keep the heat at medium and stir often; the sauce sputters as it cooks down.
    \item \textbf{Seasoning}: the adobo is salted at the blender and the \textbf{lamb} is salted on its own, so taste after frying the adobo and again after the braising juices are stirred in. Adjust gradually---it is easier to add than to correct.
    \item \textbf{Consommé~consistency}: aim for a brothy, spoonable liquid you would happily sip, not a thick sauce. Add reserved \textbf{chile~water} a little at a time; over-thinning dulls the flavor.
    \item \textbf{Griddling}: medium-high, not searing hot, or the tortilla scorches before the cheese melts. Lift an edge with the spatula to check for a crisp, deep-red underside before flipping.
    \item \textbf{Shredding}: keep the strands coarse so they hold texture in the taco and do not turn to paste when griddled.
\end{itemize}

\subsection*{Make Ahead \& Storage}
\begin{itemize}
    \item The full braise improves overnight. Cool the \textbf{shredded~lamb} and consommé, refrigerate separately up to \textit{4~days}, and skim the solidified \textbf{chile~fat} from the chilled consommé to reuse for frying.
    \item Freeze \textbf{shredded~lamb} in consommé up to \textit{3~months}; thaw overnight in the refrigerator.
    \item Griddle the tacos to order---reheated quesabirria loses its crisp edge.
    \item Reheat \textbf{lamb} and consommé gently on the stovetop; loosen the consommé with a splash of water if it has thickened.
\end{itemize}

\subsection*{Serving Suggestions}
\begin{itemize}
    \item Serve quesabirria-style: crisp folded tacos alongside a bowl of hot consommé for dipping, with extra chopped \textbf{onion}, \textbf{cilantro}, and \textbf{lime}.
    \item Offer a salsa de árbol or salsa macha for those who want more heat.
    \item Round out the meal with refried beans, Mexican rice, or a simple cabbage slaw.
    \item Beverage pairing: a Mexican lager, a michelada, or hibiscus agua fresca to cut the richness.
\end{itemize}


\end{multicols}
}

\end{document}
